\documentclass[12pt]{article}
\usepackage{amsmath}
\usepackage{amssymb}  % Provides \mathbb and other symbols
\usepackage{geometry}
\usepackage{graphicx}
\usepackage{subcaption}
\usepackage{enumitem}
\usepackage{listings}
\usepackage{hyperref}
\usepackage{ulem}
\usepackage{xcolor}
\usepackage{amsthm}
\usepackage{tcolorbox}  % For colored boxes
\usepackage[table]{xcolor} % For coloring table cells
\usepackage{array}         % For better table alignment
\usepackage{mdframed}
\usepackage{cancel}
\usepackage{amsmath,amssymb}

% Set the top margin
\geometry{top=1in} % You can adjust the measurement as needed


\begin{document}

\begin{titlepage}
\centering
\vspace*{\fill} % Adds vertical space to center the title block vertically

{\Huge Theoretical Physics for Mathematicians} \\[0.5cm] % Title 2
{\Huge Home Work 1}\\[1.5cm] % Title 3
{\Large Student: 26-712-224}\\[0.5cm] % Student ID
{\Large \today} % Date, \today adds today's date

\vspace*{\fill}
\end{titlepage}
\newpage
\noindent

\noindent

\section*{Newtonian Space-Time and the Galilei Group}

\noindent
\textbf{Exercise 1} $(\star\ \text{The Galilei group}).$
For $a \in \mathbb{R}$, $\vec{b},\vec{v} \in \mathbb{R}^3$ and
$R \in SO(3)$ define the map
\begin{align*}
g_{(a,\vec{b},\vec{v},R)}:
\mathbb{R} \times \mathbb{R}^3
&\longrightarrow
\mathbb{R} \times \mathbb{R}^3,
&
(t,\vec{x})
&\longmapsto
(t+a,R\vec{x}+\vec{v}t+\vec{b}).
\end{align*}

\begin{enumerate}
\item[(a)]
Show that the composition of two such maps is again of this form, and derive the
composition law
\begin{align*}
g_{(a_2,\vec{b}_2,\vec{v}_2,R_2)}
\circ
g_{(a_1,\vec{b}_1,\vec{v}_1,R_1)}
=
g_{(a_1+a_2,\,
R_2\vec{b}_1+a_1\vec{v}_2+\vec{b}_2,\,
R_2\vec{v}_1+\vec{v}_2,\,
R_2R_1)}.
\end{align*}
Deduce that the set $\mathcal{G}$ of all these maps is a group, compute the inverse of
$g_{(a,\vec{b},\vec{v},R)}$, and determine $\dim \mathcal{G}$.

\item[(b)]
Let
\begin{align*}
N
=
\left\{
g_{(a,\vec{b},0,I)}
\mid
a \in \mathbb{R},\ \vec{b} \in \mathbb{R}^3
\right\}
\end{align*}
be the subgroup of space-time translations. Show that $N$ is normal in
$\mathcal{G}$ and that
$\mathcal{G}/N \cong SO(3)\ltimes\mathbb{R}^3$.
Conclude that
$\mathcal{G}\cong\mathbb{R}^4\rtimes\left(SO(3)\ltimes\mathbb{R}^3\right)$.

\item[(c)]
Show that the boosts
\begin{align*}
B
=
\left\{
g_{(0,0,\vec{v},I)}
\mid
\vec{v} \in \mathbb{R}^3
\right\}
\end{align*}
form an abelian subgroup isomorphic to $(\mathbb{R}^3,+)$, but that $B$ is
\emph{not} normal in $\mathcal{G}$.

\item[(d)]
Let $\vec{x}:I\to\mathbb{E}^3$ be a trajectory and let
$\widetilde{\vec{x}}$ be its image under $g_{(a,\vec{b},\vec{v},R)}$.
Show that
\begin{align*}
\ddot{\vec{x}}\equiv 0
\quad\text{if and only if}\quad
\ddot{\widetilde{\vec{x}}}\equiv 0.
\end{align*}
Hence the notion of an inertial system is well defined on the whole group
$\mathcal{G}$.
\end{enumerate}

\newpage
\noindent
ANSWER:\\
\\
\noindent\textbf{(a)} To show that the Galilean transformations are closed under composition.\\
Take
\begin{align*}
g_{(a_1,\hat{b}_1,\hat{v}_1,R_1)}(t,\bar{x})
&=
(t+a_1,R_1\bar{x}+\hat{v}_1t+\hat{b}_1),\\
g_{(a_2,\hat{b}_2,\hat{v}_2,R_2)}(t,\bar{x})
&=
(t+a_2,R_2\bar{x}+\hat{v}_2t+\hat{b}_2).
\end{align*}
Then
\begin{align*}
g_{(a_2,\hat{b}_2,\hat{v}_2,R_2)}
\circ
g_{(a_1,\hat{b}_1,\hat{v}_1,R_1)}(t,\bar{x})
&=
g_{(a_2,\hat{b}_2,\hat{v}_2,R_2)}
(t+a_1,R_1\bar{x}+\hat{v}_1t+\hat{b}_1)\\
&=
\left(
t+a_1+a_2,
R_2[R_1\bar{x}+\hat{v}_1t+\hat{b}_1]
+\hat{v}_2(t+a_1)+\hat{b}_2
\right)\\
&=
\left(
t+a_1+a_2,
(R_2R_1)x
+(R_2\hat{v}_1)t
+R_2\hat{b}_1
+\hat{v}_2t
+\hat{v}_2a_1
+\hat{b}_2
\right)\\
&=
\left(
t+a_1+a_2,
(R_2R_1)\bar{x}
+(R_2\hat{v}_1+\hat{v}_2)t
+(R_2\hat{b}_1+\hat{v}_2a_1+\hat{b}_2)
\right)\\
&=
g_{(a_1+a_2,\,
R_2\hat{b}_1+\hat{v}_2a_1+\hat{b}_2,\,
R_2\hat{v}_1+\hat{v}_2,\,
R_2R_1)}.
\tag{*}
\end{align*}
Now $R_1,R_2\in SO(3)\Rightarrow R_2R_1\in SO(3)$.  Also: $R_2\bar{v_1}+\bar{v_2}\in\mathbb{R}^{3}$, $R_2\bar{b_1}+\bar{v_2}a+\bar{b_2} \in\mathbb{R}^{3}$.\\
Hence $(*)$ is also a Galilean transformation.\\
\\
The identity element of this group will leave $(t,\bar{x})$ fixed. Thus:
\begin{align*}
g_{(a,\bar{b},\bar{v},R)}(t,\bar{x})
&=(t+a,R\bar{x}+\bar{v}t+\bar{b})\\
&=(t,\bar{x}).
\end{align*}
It is clear that: $a=0$, $\bar{v}=\bar{0}$, $\bar{b}=\bar{0}$ and $R=I$.\\
Hence $g_{(0,\bar{0},\bar{0},I)}$ is the identity element for the group $\mathcal{G}$.\\
\\
Now lets find the inverse transformation for
$g_{(a,\bar{b},\bar{v},R)}\in G$.  Such that:
\begin{align*}
g_{(c,\bar{d},\bar{u},T)}
\circ g_{(a,\bar{b},\bar{v},R)}
&=g_{(0,\bar{0},\bar{0},I)}. \tag{1}\\
g_{(a,\bar{b},\bar{v},R)}
\circ g_{(c,\bar{d},\bar{u},T)}
&=g_{(0,\bar{0},\bar{0},I)}. \tag{2}
\end{align*}
From $(*)$ above we know that
\begin{align*}
g_{(c,\bar{d},\bar{u},T)}
\circ g_{(a,\bar{b},\bar{v},R)}
&=
g_{(a+c,\,
T\bar{b}+\bar{u}a+\bar{d},\,
T\bar{v}+\bar{u},\,
TR)}\\
&=
g_{(0,\bar{0},\bar{0},I)}.
\end{align*}
Clearly:
\begin{align*}
a+c&=0
\quad\Rightarrow\quad
c=-a,\\
TR&=I
\quad\Rightarrow\quad
T=R^{-1}
\qquad
(\text{possible as }R\in SO(3)),\\
T\bar{v}+\bar{u}&=\bar{0}
\quad\Rightarrow\quad
\bar{u}=-T\bar{v}=-R^{-1}\bar{v},\\
T\bar{b}+\bar{u}a+\bar{d}&=\bar{0}
\quad\Rightarrow\quad
R^{-1}\bar{b}-(R^{-1}\bar{v})a+\bar{d}=\bar{0}\\
&\quad\Rightarrow\quad
\bar{d}=R^{-1}(\bar{v}a-\bar{b}).
\end{align*}
The same result applies if we take $(2)$. Hence
\begin{align*}
g_{(a,\bar{b},\bar{v},R)}^{-1}
=
g_{(-a,\,
R^{-1}(\bar{v}a-\bar{b}),\,
-R^{-1}\bar{v},\,
R^{-1})}.
\end{align*}
The above three properties: closure, inverse and identity show that $\mathcal{G}$, the set of Galilean transformations, is indeed a group.\\
The total number of dimensions for $\mathcal{G}$ is 10 = 1 (time) + 3 (position) + 3 (volocity) + 3 (rotation matrix).\\
\\
\noindent\textbf{(b)}\\
I didn't know what a "semidirect product" was - I understand it as follows:\\
\begin{tcolorbox}[colframe=red,coltext=red]
A \textbf{semidirect product} is a group $G$ built from subgroups $N$ and $H$ such that
\begin{align*}
N \trianglelefteq G,\qquad N\cap H=\{e\},\qquad G=NH.
\end{align*}
Unlike a direct product, $N$ and $H$ need not commute. One subgroup acts on the other by conjugation.
\end{tcolorbox}
\noindent
To demonstrate that: $N=\{g_{(a, \bar{b}, \bar{0}, I)} \; | \; a \in \mathbb{R}, \bar{b} \mathbb{R}^3  \}$ is normal in $\mathcal{G}$ we need to prove that:
given any $g \in \mathcal{G}$ and $n \in N$, that:
$$ g \circ n \circ g^{-1} \in N $$
This is just subscript manipulation using equation (*) repeatedly.  Take $n=g_{(a, \bar{b}, \bar{0}, I)}$ and $g=g_{(a_2, \bar{b_2}, \bar{v_2}, R_2)}$.
Then:
\begin{align*}
  g_{(a_2, \bar{b_2}, \bar{v_2}, R_2)} \circ n \circ g_{(a_2, \bar{b_2}, \bar{v_2}, R_2)}^{-1} &=  g_{(a_2, \bar{b_2}, \bar{v_2}, R_2)} \circ g_{(a, \bar{b}, \bar{0}, I)} \circ g_{(-a_2, R_2^{-1}(\bar{v_2} a_2 - \bar{b_2}), -R_2^{-1}\bar{v_2}, R_2^{-1})} \\
  &= g_{(a_2 + a, R_2\bar{b} + a \bar{v_2}+ \bar{b_2}, \bar{v_2}, R_2 I )} \circ g_{(-a_2, R_2^{-1}(\bar{v_2} a_2 - \bar{b_2}), -R_2^{-1}\bar{v_2}, R_2^{-1})} \\
  &= g_{(a_2 + a, R_2\bar{b} + a \bar{v_2}+ \bar{b_2}, \bar{v_2}, R_2 )} \circ g_{(-a_2, R_2^{-1}(\bar{v_2} a_2 - \bar{b_2}), -R_2^{-1}\bar{v_2}, R_2^{-1})} \\
  &= g_{(a,R_2\bar{b} + \bar{v_2} a,\bar{0},I)}
\end{align*}
Which is again an element of $N$.  Hence $N$ is a normal subgroup of $\mathcal{G}$.\\
\\
Clearly $N$ has dimension 4 (1-time and 3-position vector) - which leaves 10-4=6 dimensions to be accounted for in $\mathcal{G}/N$.  
The remaining 6 dimensions are accounted for by the velocity and the rotations - which are independent, so:  $SO(3) \ltimes \mathbb{R}^3$
and hence: $\mathcal{G}/N \cong SO(3) \ltimes \mathbb{R}^3$.  It would be a simple procedure to show that there is an isomorphism between
$\kappa : (\mathcal{G}/N, SO(3) \ltimes \mathbb{R}^3) \to \mathcal{G}$.\\
\\
\textbf{(c)}\\
We will show that $B = \{g_{(0, \bar{0}, \bar{v}, I)} \; | \; \bar{v} \in  \mathbb{R}^3\}$ is an abelian group.\\
This is simply demonstrated as follows (and using (*)):
\begin{align*}
  g_{(0, \bar{0}, \bar{v_1}, I)} \circ g_{(0, \bar{0}, \bar{v_2}, I)} &= g_{(0, \bar{0}, \bar{v_2}, I)} \circ g_{(0, \bar{0}, \bar{v_1}, I)}\tag{**} \\
  &= g_{(0, \bar{0}, \bar{v_1} + \bar{v_2}, I)} \in B 
\end{align*}
This shows that composition is both commutative and closed on the set $B$.  The inverse of $g_{(0, \bar{0}, \bar{v}, I)}$ is obviously: $g_{(0, \bar{0}, -\bar{v}, I)}$.
Hence $B$ is an abelian group.\\
\\
We can clearly define an isomorphism between $(\mathbb{R}^3, +)$ and $(B, \circ)$ - which is just a mapping between velocities and boosts as follows:
\begin{align*}
  \phi : \mathbb{R}^3 \to B \text{  where:  } \phi(\bar{v}) = g_{(0, \bar{0}, \bar{v}, I)}
\end{align*}
From which:
\begin{align*}
  \phi(v_1 + v_2) =  g_{(0, \bar{0}, \bar{v_1} + \bar{v_2}, I)} = g_{(0, \bar{0}, \bar{v_1}, I)} \circ g_{(0, \bar{0}, \bar{v_2}, I)} = \phi(v_1) \circ \phi(v_2)
\end{align*}
\\
Lastly, we must show that $B \not\trianglelefteq \mathcal{G}$. \\
We will proceed as we did in (b).  Take any $g=g_{(a, \bar{b}, \bar{u},R)} \in \mathcal{G}$ and $b=g_{(0, \bar{0}, \bar{v}, I)} \in B$; then:
\begin{align*}
   g \circ b \circ g^{-1} &= g_{(a, \bar{b}, \bar{u},R)} \circ g_{(0, \bar{0}, \bar{v}, I)} \circ g_{(-a, R^{-1}(\bar{u}a) - \bar{b}, -R^{-1}\bar{u},R^{-1})} \\
   &= g_{(a, \bar{b}, R\bar{v} + \bar{u},R)} \circ g_{(-a, R^{-1}(\bar{u}a) - \bar{b}, -R^{-1}\bar{u},R^{-1})} \\
   &= g_{(0, -a R \bar{v}, R\bar{v},I)} \notin B
\end{align*}
So, $B$ is NOT a normal subgroup of $\mathcal{G}$. \\
\\
\textbf{(d)}\\
For the given transformation we have:
\begin{align*}
t' &= t+a \tag{3}\\
\vec{x}' &= R\vec{x}+\vec{v}t+\vec{b} \tag{4}
\end{align*}
$(3)\Rightarrow \Delta t=\Delta t'$ as time difference between any two events in the two different frames.
Hence:
\begin{align*}
\frac{d}{dt'}=\frac{d}{dt}.
\end{align*}
If we then differentiate (4) w.r.t time we get:
\begin{align*}
\dot{\vec{x}}'=\dot{R}\vec{x}+R\dot{\vec{x}}+\vec{v}
\end{align*}
Again:
\begin{align*}
\ddot{\vec{x}}'=\ddot{R}\vec{x}+\dot{R}\dot{\vec{x}}+\dot{R}\dot{\vec{x}}+R\ddot{\vec{x}}
\end{align*}
For a Galilean transformation $\dot{R}=0$, otherwise
we are not transforming between inertial frames. \\
Hence:
\begin{align*}
\ddot{\vec{x}}'=R\ddot{\vec{x}} \tag{5}
\end{align*}
If $\ddot{\vec{x}}'=\vec{0}\Rightarrow\ddot{\vec{x}}=\vec{0}$, since $R$ is not singular.
Also $\ddot{\vec{x}}=\vec{0}\Rightarrow\ddot{\vec{x}}'=0$ follows directly from (5).  This means that 
Galilean transformations are valid transformations between inertial reference frames.
\newpage
\noindent
\textbf{Exercise 4} ($\star$ Conservation laws for systems of point masses). Consider $N$ point masses
$m_j$ with trajectories $\vec{x}_j(t)\in\mathbb{E}^3$ obeying
\begin{align*}
m_j\ddot{\vec{x}}_j=\vec{k}_j=\sum_{i\neq j}\vec{k}_{ij}+\vec{K}_j,
\qquad j=1,\ldots,N,
\end{align*}
where $\vec{k}_{ij}$ is the internal force exerted on $j$ by $i$ and $\vec{K}_j$ is external.\\
Assume Newton's
third law $\vec{k}_{ij}=-\vec{k}_{ji}$. Write $\vec{p}_j=m_j\dot{\vec{x}}_j$, $\vec{P}=\sum_j\vec{p}_j$,
$M=\sum_jm_j$, $\vec{X}=M^{-1}\sum_jm_j\vec{x}_j$ and
\begin{align*}
\vec{J}=\sum_j\vec{x}_j\wedge\vec{p}_j.
\end{align*}

\textbf{(a)} Show that $\dot{\vec{P}}=\sum_j\vec{K}_j$, and deduce that $\vec{P}$ is conserved for a closed system.

\textbf{(b)} Deduce the centre-of-mass theorem: for a closed system
\begin{align*}
\vec{X}(t)=\vec{X}_0+\vec{V}t.
\end{align*}

\textbf{(c)} Show that
\begin{align*}
\dot{\vec{J}}=\sum_j\vec{x}_j\wedge\vec{K}_j
\end{align*}
\hspace*{1cm}provided the internal forces are central, i.e.
\begin{align*}
\vec{k}_{ji}
=
\frac{\vec{x}_j-\vec{x}_i}{|\vec{x}_j-\vec{x}_i|}
f_{ji}(|\vec{x}_j-\vec{x}_i|).
\end{align*}
\hspace*{1cm}Where exactly does centrality enter? Give an example of forces obeying \\
\hspace*{2cm}$\vec{k}_{ji}=-\vec{k}_{ij}$ for which $\vec{J}$ is not conserved.\\

\textbf{(d)} (Two-body reduction.) Let $N=2$ with central internal forces and no external\\
\hspace*{1cm}force.
Set $\vec{r}=\vec{x}_1-\vec{x}_2$ and $\mu=\frac{m_1m_2}{M}$. Show that

\begin{align*}
\mu\ddot{\vec{r}}
=
\frac{\vec{r}}{|\vec{r}|}f_{12}(|\vec{r}|),
\end{align*}

\hspace*{1cm}and prove K\"onig's decompositions

\begin{align*}
T
&=
\frac{1}{2}M|\dot{\vec{X}}|^2
+
\frac{1}{2}\mu|\dot{\vec{r}}|^2,
\\
\vec{J}
&=
M\vec{X}\wedge\dot{\vec{X}}
+
\mu\vec{r}\wedge\dot{\vec{r}}.
\end{align*}

\newpage
\noindent
ANSWER:\\
\\
\textbf{(a)}
\begin{align*}
\vec{P} &= \sum_j \vec{p}_j \\
\Rightarrow \dot{\vec{P}}
&= \sum_j \dot{\vec{p}}_j \\
&= \sum_j m_j \ddot{\vec{x}}_j
\qquad \qquad \text{(as $m_j$ is constants)} \\
&= \sum_j \left\{
\left[\sum_{i\neq j}\vec{k}_{ji}\right]+\vec{K}_j
\right\}. \tag{4.1}
\end{align*}

Now if we take $\vec{k}_{ii}=0\;(\forall i)$ then (4.1) becomes,
\begin{align*}
\dot{\vec{P}}
=
\sum_j\left\{
\sum_i \vec{k}_{ji}+\vec{K}_j
\right\}.
\end{align*}
\hspace*{1cm}Since: $\vec{k}_{ij}=-\vec{k}_{ji} \Rightarrow \sum_j\sum_i\vec{k}_{ij}=\vec{0}$.
\begin{align*}
\therefore\quad \dot{\vec{P}}=\sum_j\vec{K}_j.
\end{align*}
\hspace*{1cm}For a closed system each $\vec{K}_j=\vec{0}$
\begin{align*}
&\Rightarrow\dot{\vec{P}}=\vec{0} \\
&\Rightarrow\vec{P}\text{ is constant}.
\end{align*}
\vspace{1cm}
\noindent
\textbf{(b)}
\begin{align*}
M\vec{X}
&= M\left\{\frac{1}{M}\sum_j m_j\vec{x}_j\right\} \\
&= \sum_j m_j\vec{x}_j
\end{align*}
Differentiating w.r.t. time:
\begin{align*}
M\dot{\vec{X}}
&= \sum_j m_j\dot{\vec{x}}_j \\
&= \vec{P} \qquad\qquad (\text{a constant}). \\
&= M\vec{V} \qquad \;\;\;\; (\text{where }\vec{V}\text{ is a constant}).
\end{align*}
\begin{align*}
\Rightarrow\quad \dot{\vec{X}} &= \vec{V} \\
\Rightarrow\quad \vec{X} &= \vec{V}t+\vec{X}_0
\qquad (X_0=\text{const. of integration}).
\end{align*}
\noindent
\textbf{(c)}
\begin{align*}
\vec{J} &= \sum_j \vec{x}_j \wedge \vec{p}_j \\
\Rightarrow \dot{\vec{J}} &= \sum_j \dot{\vec{x}}_j \wedge \vec{p}_j + \sum_j \vec{x}_j \wedge \dot{\vec{p}}_j \tag{4.2}
\end{align*}
Now $\vec{p}_j=m_j\dot{\vec{x}}_j \Rightarrow \dot{\vec{x}}_j\wedge\vec{p}_j=0\quad\forall j$. Hence $(4.2)$ becomes:
\begin{align*}
\dot{\vec{J}}
&= \sum_j \vec{x}_j \wedge \dot{\vec{p}}_j \\
&= \sum_j \vec{x}_j \wedge \vec{k}_j \\
&= \sum_j \vec{x}_j \wedge
\left\{
\sum_{i\neq j}\vec{k}_{ji}+\vec{K}_j \tag{4.3}
\right\}.
\end{align*}
Now
\begin{align*}
\sum_j \vec{x}_j \wedge \sum_{i\neq j}\vec{k}_{ji}=0
\end{align*}
since, for any $i\neq j$ (using centrality here), force is along line joining the bodies:
\begin{align*}
\vec{x}_j\wedge\vec{k}_{ji}
&=
\vec{x}_j\wedge
\frac{\vec{x}_j-\vec{x}_i}{|\vec{x}_j-\vec{x}_i|}
f_{ji}(|\vec{x}_j-\vec{x}_i|), \\
\vec{x}_i\wedge\vec{k}_{ij}
&=
\vec{x}_i\wedge
\frac{\vec{x}_i-\vec{x}_j}{|\vec{x}_i-\vec{x}_j|}
f_{ij}(|\vec{x}_i-\vec{x}_j|).
\end{align*}
Their sum is
\begin{align*}
(\vec{x}_j\wedge\vec{k}_{ji})+(\vec{x}_i\wedge\vec{k}_{ij})
&=
(\vec{x}_j-\vec{x}_i)\wedge
\frac{\vec{x}_j-\vec{x}_i}{|\vec{x}_j-\vec{x}_i|}
f_{ij}(|\vec{x}_j-\vec{x}_i|) \\
&=0
\qquad \text{as } f_{ij}=f_{ji}\quad(\text{Newton's 3rd}).
\end{align*}
Thus $(4.3)$ becomes:
\begin{align*}
\dot{\vec{J}}=\sum_j \vec{x}_j\wedge\vec{K}_j.
\end{align*}

\newpage
\noindent
\textbf{(d)}\\
(i)
\begin{align*}
\bar{k}_1 &= m_1\ddot{\bar{x}}_1
= \frac{\bar{x}_1-\bar{x}_2}{|\bar{x}_1-\bar{x}_2|}
f_{12}(|\bar{x}_1-\bar{x}_2|) \\
\bar{k}_2 &= m_2\ddot{\bar{x}}_2
= \frac{\bar{x}_2-\bar{x}_1}{|\bar{x}_1-\bar{x}_2|}
f_{21}(|\bar{x}_2-\bar{x}_1|).
\end{align*}
\begin{align*}
\Rightarrow\quad
\ddot{\bar{x}}_1
&= \frac{1}{m_1}
\frac{\bar{x}_1-\bar{x}_2}{|\bar{x}_1-\bar{x}_2|}
f_{12}(|\bar{x}_1-\bar{x}_2|) \\
\ddot{\bar{x}}_2
&= \frac{1}{m_2}
\frac{\bar{x}_2-\bar{x}_1}{|\bar{x}_2-\bar{x}_1|}
f_{21}(|\bar{x}_2-\bar{x}_1|).
\end{align*}
\begin{align*}
\Rightarrow\quad
\ddot{\bar{x}}_2-\ddot{\bar{x}}_1
&=
\frac{1}{m_2}
\frac{\bar{x}_2-\bar{x}_1}{|\bar{x}_2-\bar{x}_1|}
f_{21}(|\bar{x}_2-\bar{x}_1|)
+
\frac{1}{m_1}
\frac{\bar{x}_2-\bar{x}_1}{|\bar{x}_2-\bar{x}_1|}
f_{12}(|\bar{x}_2-\bar{x}_1|) \\
&=
\frac{m_1+m_2}{m_1m_2}
\frac{\bar{x}_2-\bar{x}_1}{|\bar{x}_2-\bar{x}_1|}
f_{12}(|\bar{x}_2-\bar{x}_1|).
\end{align*}
\begin{align*}
\Rightarrow\quad
\frac{m_1m_2}{m_1+m_2}
(\ddot{\bar{x}}_2-\ddot{\bar{x}}_1)
=
\frac{\bar{x}_2-\bar{x}_1}{|\bar{x}_2-\bar{x}_1|}
f_{12}(|\bar{x}_2-\bar{x}_1|)
\end{align*}
i.e.
\begin{align*}
\mu\ddot{\bar{r}}
=
\frac{\bar{r}}{|\bar{r}|}f_{12}(|\bar{r}|).
\end{align*}
It should be noted that $f_{12}(|r|) = f_{21}(|r|)$ must be the case for Newton's 3rd law to hold.\\
\\
\textbf{(ii)}
\begin{align*}
T &= \frac{1}{2}m_1\dot{\bar{x}}_1^2 + \frac{1}{2}m_2\dot{\bar{x}}_2^2
\end{align*}
Now
\begin{align*}
\bar{X}
&=\frac{1}{m_1+m_2}\left[m_1\bar{x}_1+m_2\bar{x}_2\right] \\
\Rightarrow\quad
\dot{\bar{X}}
&=\frac{1}{m_1+m_2}\left[m_1\dot{\bar{x}}_1+m_2\dot{\bar{x}}_2\right] \\
\Rightarrow\quad
|\dot{\bar{X}}|^2
&=\dot{\bar{X}}^2
=\frac{1}{(m_1+m_2)^2}
\left[m_1^2\dot{\bar{x}}_1^2
+2m_1m_2\dot{\bar{x}}_1\cdot\dot{\bar{x}}_2
+m_2^2\dot{\bar{x}}_2^2\right] \\
\Rightarrow\quad
\frac{1}{2}(m_1+m_2)\dot{\bar{X}}^2
&=\frac{1}{2(m_1+m_2)}
\left[m_1^2\dot{\bar{x}}_1^2
+2m_1m_2\dot{\bar{x}}_1\cdot\dot{\bar{x}}_2
+m_2^2\dot{\bar{x}}_2^2\right]
\end{align*}
Hence, using $M=(m_1+m_2)$, we obtain:
\begin{align*}
T-\frac{1}{2}M\dot{\bar{X}}^2
&=
\left[
\frac{1}{2}m_1\dot{\bar{x}}_1^2
+\frac{1}{2}m_2\dot{\bar{x}}_2^2
\right]
-\frac{m_1^2}{2(m_1+m_2)}\dot{\bar{x}}_1^2
-\frac{m_1m_2}{m_1+m_2}\dot{\bar{x}}_1\cdot\dot{\bar{x}}_2 -\frac{m_2^2}{2(m_1+m_2)}\dot{\bar{x}}_2^2 \\
&=
\frac{1}{2}\dot{\bar{x}}_1^2
\left[m_1-\frac{m_1^2}{m_1+m_2}\right]
+\frac{1}{2}\dot{\bar{x}}_2^2
\left[m_2-\frac{m_2^2}{m_1+m_2}\right] -\frac{m_1m_2}{m_1+m_2}
\dot{\bar{x}}_1\cdot\dot{\bar{x}}_2 \\
&=
\frac{1}{2}\frac{m_1m_2}{m_1+m_2}
\left[
\dot{\bar{x}}_1^2
-2\dot{\bar{x}}_1\cdot\dot{\bar{x}}_2
+\dot{\bar{x}}_2^2
\right] \\
&=
\frac{1}{2}\frac{m_1m_2}{m_1+m_2}
\left(\dot{\bar{x}}_2-\dot{\bar{x}}_1\right)^2 \\
&=
\frac{1}{2}\mu\dot{\bar{r}}^2
\end{align*}
i.e
\begin{align*}
T=\frac{1}{2}M\dot{\bar{X}}^2+\frac{1}{2}\mu\dot{\bar{r}}^2
\end{align*}
\\
(iii) \begin{align*}
\bar{J} &= \bar{x}_1 \wedge m_1\dot{\bar{x}}_1 + \bar{x}_2 \wedge m_2\dot{\bar{x}}_2
\end{align*}
Now
\begin{align*}
M\bar{X}\wedge\dot{\bar{X}}
&= (m_1+m_2)\frac{1}{m_1+m_2}(m_1\bar{x}_1+m_2\bar{x}_2)
\wedge\frac{1}{m_1+m_2}(m_1\dot{\bar{x}}_1+m_2\dot{\bar{x}}_2) \\
&= \frac{1}{m_1+m_2}(m_1\bar{x}_1+m_2\bar{x}_2)
\wedge(m_1\dot{\bar{x}}_1+m_2\dot{\bar{x}}_2) \\
&= \frac{1}{m_1+m_2}\left[
m_1^2\bar{x}_1\wedge\dot{\bar{x}}_1
+m_2^2\bar{x}_2\wedge\dot{\bar{x}}_2
+m_1m_2\bar{x}_1\wedge\dot{\bar{x}}_2
+m_1m_2\bar{x}_2\wedge\dot{\bar{x}}_1
\right]
\end{align*}
Hence:
\begin{align*}
\bar{J}-M\bar{X}\wedge\dot{\bar{X}}
&=
\left[
m_1\bar{x}_1\wedge\dot{\bar{x}}_1
+m_2\bar{x}_2\wedge\dot{\bar{x}}_2
\right] \\
&\quad-\frac{1}{m_1+m_2}\left[
m_1^2\bar{x}_1\wedge\dot{\bar{x}}_1
+m_2^2\bar{x}_2\wedge\dot{\bar{x}}_2
+m_1m_2\bar{x}_1\wedge\dot{\bar{x}}_2
+m_1m_2\bar{x}_2\wedge\dot{\bar{x}}_1
\right] \\
&=
\frac{m_1m_2}{m_1+m_2}\left[
\bar{x}_1\wedge\dot{\bar{x}}_1
+\bar{x}_2\wedge\dot{\bar{x}}_2
-\bar{x}_1\wedge\dot{\bar{x}}_2
-\bar{x}_2\wedge\dot{\bar{x}}_1
\right] \\
&=
\frac{m_1m_2}{m_1+m_2}
(\bar{x}_2-\bar{x}_1)\wedge
(\dot{\bar{x}}_2-\dot{\bar{x}}_1) \\
&= \mu\bar{r}\wedge\dot{\bar{r}}
\end{align*}
Hence:
\begin{align*}
\bar{J}=M\bar{X}\wedge\dot{\bar{X}}+\mu\bar{r}\wedge\dot{\bar{r}}
\end{align*}

\end{document}